%%%% CAPÍTULO 6 -- CONSIDERAÇÕES FINAIS
%%
%% Este capítulo consolida a discussão de limitações e trabalhos
%% futuros do estudo. Trechos hoje dispersos em
%% `cap-experimentos-resultados.tex` (subsec:IRISMobileLimitacoes,
%% subsec:DeployLimitacoes e parágrafos isolados da subsec:NPITrocaDados)
%% devem migrar para cá; em seus pontos de origem deve permanecer
%% apenas uma frase de remissão a este capítulo.

\chapter{Considerações finais}\label{cap:consideracoesFinais}

% ======================================================================
\section{Síntese do trabalho}\label{sec:sinteseTrabalho}
% ======================================================================

O presente trabalho propôs o Padrão de Rótulos e Interfaces para
Sistemas Médicos (PRISM), um conjunto de diretrizes e processos de
interoperabilidade orientado à pesquisa clínico-hospitalar baseada em
biossinais. A motivação, apresentada no Capítulo~\ref{cap:introducao},
parte do diagnóstico de que a fragmentação dos dados produzidos em
projetos de pesquisa biomédica --- hoje confinados a silos próprios e
formatos não normalizados --- compromete a reprodutibilidade dos
estudos, encarece a integração entre instituições e prejudica a
qualidade dos modelos baseados em aprendizado de máquina aplicados ao
contexto clínico.

A proposta foi materializada por meio do Nó de Pesquisa Interoperável
(NPI), do \textit{Interoperable Research Interface System} (IRIS), de
sua versão móvel especializada em captura (IRIS \textit{Mobile}) e da
adaptação do \textit{firmware} do NeuroEstimulator, descritos em
detalhe no Capítulo~\ref{cap:desenvolvimento}. A operação fim a fim
desses artefatos foi exercitada por meio de dois testes
complementares no Capítulo~\ref{cap:resultados}: a validação
quantitativa da cadeia de aquisição com gerador de funções e a
demonstração da sincronização entre dois NPIs federados.

Os objetivos específicos declarados na
Subseção~\ref{subsec:objetivosEspecificos} foram atendidos da
seguinte forma. A revisão da literatura e do estado da arte foi
conduzida no Capítulo~\ref{cap:referencialTeorico} e na
Seção~\ref{sec:estadoDaArte}, contemplando iniciativas
governamentais (RNDS, USCDI, EHDS, NHS, \textit{Shared
Pan-Canadian} e \textit{My Health Record}) e trabalhos acadêmicos em
torno de perfis FHIR e ontologias clínicas. O desenvolvimento do NPI
foi concretizado na Seção~\ref{sec:DesenvolvimentoNPI}, com sua
arquitetura interna, modelo de dados, \textit{middleware} de
identificação e protocolos de troca documentados em subseções
próprias. A modificação, adaptação e integração do projeto de
captura de eletromiografia de superfície foi tratada na
Seção~\ref{sec:adaptacaoNeuroEstimulator}, complementada pela
construção do IRIS \textit{Mobile} na
Seção~\ref{sec:desenvolvimentoIRISMobile}. O teste do resultado
final da integração foi conduzido nas Seções
\ref{sec:resultadosCadeiaAquisicao} e
\ref{sec:resultadosSincronizacao}, e a avaliação dos resultados
obtidos foi sintetizada no Quadro~\ref{quad:resultadosHipoteses} e na
discussão de limitações e trabalhos futuros que se segue.

% ======================================================================
\section{Limitações do estudo}\label{sec:limitacoesEstudo}
% ======================================================================

Esta seção consolida as limitações reconhecidas ao longo do trabalho.
A consolidação aqui é deliberada: optou-se por recolher as restrições
apontadas pontualmente nos capítulos anteriores em uma análise única,
agrupada por categoria, de modo a permitir ao leitor avaliar o
impacto agregado das escolhas feitas pelo autor sobre a abrangência
da validação experimental.

% ----------------------------------------------------------------------
\subsection{Limitações de escopo experimental}\label{subsec:limitacoesEscopo}
% ----------------------------------------------------------------------

O recorte experimental adotado é deliberadamente estreito. O
trabalho restringe-se a uma única modalidade de biossinal (sEMG), a
um único projeto-exemplo e a uma federação composta por apenas dois
nós, conforme já justificado na visão geral do
Capítulo~\ref{cap:desenvolvimento}. Essas escolhas, decorrentes das
restrições de equipe (um único autor) e de prazo (ciclo acadêmico de
um ano), preservam a capacidade do experimento de testar a hipótese
central, mas limitam a generalização imediata dos resultados a
ambientes de operação produtiva.

A topologia em uma única máquina virtual Azure
(Subseção~\ref{subsec:DeployAzure}) restringe a abrangência da
validação aos aspectos lógicos da sincronização entre nós --- o
contrato de mensagens, a sequência das fases do \textit{handshake}, a
idempotência e a resolução de conflitos ---, mas não captura efeitos
típicos de ambientes geograficamente distribuídos, como latência
interinstitucional, partições de rede e divergência de relógios entre
nós.

A coleta de sinais para alimentar a sincronização foi realizada a
partir de um único voluntário (o próprio autor), em consonância com
o foco da pesquisa na estrutura dos dados e em sua sincronização
entre nós, e não na qualidade do sinal capturado ou na variabilidade
inter-sujeitos. Os procedimentos de aquisição foram aprovados pelo
Comitê de Ética em Pesquisa com Seres Humanos da UTFPR sob parecer
nº~7.218.395, conforme detalhado no Capítulo~\ref{cap:metodologia}.

% ----------------------------------------------------------------------
\subsection{Limitações do dispositivo (NeuroEstimulator)}\label{subsec:limitacoesDispositivo}
% ----------------------------------------------------------------------

A escolha pelo NeuroEstimulator como dispositivo-exemplo trouxe a
restrição de não modificar o \textit{hardware} existente (requisito 4
da Subseção~\ref{subsec:NERequisitos}), o que limita o universo de
soluções possíveis para o estágio de aquisição. Em particular,
inviabiliza a substituição do ADS1115 por um conversor de taxa mais
alta e impede a troca do módulo Bluetooth Classic HC-06 por um
\textit{Bluetooth Low Energy} (BLE) ou Wi-Fi de maior largura de
banda.

Como consequência direta, a operação do conversor a 860\,SPS, seguida
da decimação 4:1 descrita na Subseção~\ref{subsubsec:NEDecimacao},
fixa a taxa efetiva de saída em 215\,SPS e a frequência de Nyquist da
cadeia em 107,5\,Hz. Esse valor é inferior ao extremo superior da
banda fisiológica do sEMG (cerca de 500\,Hz), conforme caracterizado
quantitativamente na Subseção~\ref{subsec:resultadosCadeiaLimitacoes}.
Essa restrição é uma propriedade do dispositivo-exemplo
adotado, e não do PRISM.

Do ponto de vista de comunicação, a dependência de Bluetooth Classic
SPP transfere-se para os clientes do dispositivo, restringindo o
IRIS \textit{Mobile} à plataforma Android e impedindo, sem mediação,
o uso do aplicativo em iOS ou em ambiente \textit{web} BLE.

% ----------------------------------------------------------------------
\subsection{Limitações do IRIS Mobile}\label{subsec:limitacoesIRISMobile}
% ----------------------------------------------------------------------

%% NOTA: este conteúdo foi consolidado a partir da subsec:IRISMobileLimitacoes
%% originalmente em capitulos/cap-experimentos-resultados.tex (linhas 773-786).
%% Sugere-se remover o conteúdo daquela subseção e deixar apenas uma frase
%% remetendo a esta subseção.

O IRIS \textit{Mobile} foi mantido enxuto e focado nas necessidades do
experimento descrito no Capítulo~\ref{cap:desenvolvimento}, com cinco
limitações reconhecidas que devem ser explicitadas para orientar
trabalhos futuros.

A primeira limitação diz respeito à \textbf{resiliência da fila de
sincronização}. O \texttt{SyncService}, em sua implementação atual,
mantém em memória o contador de tentativas de cada entidade
reenfileirada por \textit{backoff} exponencial; um reinício do
aplicativo zera esse contador e reapresenta a entidade ao laço de
sincronização como se fosse a primeira tentativa. O comportamento é
benigno em uso típico, mas em cenários de uso prolongado com NPI
indisponível a evolução natural seria uma fila persistente, conforme
discutido na Subseção~\ref{subsec:trabalhosFuturosIrisMobile}.

A segunda limitação refere-se à \textbf{ausência de suporte a
múltiplos dispositivos simultâneos}. O \texttt{BluetoothContext}, por
desenho, mantém uma única conexão SPP ativa por vez e não introduz
mecanismos de \textit{ownership} ou de roteamento entre múltiplos
dispositivos. Embora o experimento atual envolva apenas um
NeuroEstimulator por sessão, cenários futuros que demandem aquisição
multicanal a partir de dispositivos distintos exigirão uma
reformulação substancial do contexto Bluetooth e do esquema de
\texttt{recordings}.

A terceira limitação atinge a \textbf{visualização do sinal em tempo
real}. O \texttt{SEMGChart} apresenta o sinal bruto em janela de
quatro segundos sem opções configuráveis de filtragem em tempo real
(passa-faixa, retificação, envoltória), de marcadores temporais
sobrepostos ao traço ou de eixos auxiliares. A ausência decorre do
compromisso entre fidelidade visual e custo de re-renderização:
introduzir filtragens adicionais em tempo de execução elevaria a
frequência de atualização exigida e comprometeria a fluidez observada
no aparelho de testes.

A quarta limitação é a \textbf{granularidade temporal das anotações}.
O vínculo entre rótulo e trecho de sinal é estabelecido pelo
\texttt{created\_at} da anotação e pelo \texttt{timestamp} dos
quadros. Embora suficiente para o experimento conduzido, esse vínculo
é insuficiente para usos clínicos que demandem precisão sub-segundo
ou rótulos temporais fechados (com início e fim explícitos sobre o
sinal).

A quinta e última limitação diz respeito à \textbf{cobertura de
plataformas}. A dependência de Bluetooth Classic restringe o
aplicativo, na prática, ao Android: a versão \textit{web} gerada pelo
\textit{Expo} opera apenas com BLE, e o iOS não expõe SPP a
aplicativos de terceiros.

% ----------------------------------------------------------------------
\subsection{Limitações operacionais e de segurança}\label{subsec:limitacoesOperacionais}
% ----------------------------------------------------------------------

%% NOTA: este conteúdo foi consolidado a partir da subsec:DeployLimitacoes
%% originalmente em capitulos/cap-experimentos-resultados.tex (linhas 984-997).
%% Sugere-se reduzir aquela subseção a um parágrafo introdutório que
%% remeta a esta.

A implantação experimental privilegia a validação dos contratos de
interoperabilidade do PRISM em detrimento das características
operacionais que diferenciam um ambiente de pesquisa de uma
instalação de produção. As principais limitações decorrentes desse
recorte são sintetizadas a seguir.

A limitação mais visível é a \textbf{coexistência de HTTP e HTTPS} sem
que a segunda variante seja, hoje, utilizável como caminho operacional.
O NGINX já provê \textit{reverse proxy} HTTPS para os quatro serviços
publicados (NPI A/B e IRIS \textit{web} A/B), mas o certificado
autoassinado oferecido (com \textit{Common Name} fixado no IP
\texttt{104.45.216.5} e sem \textit{Subject Alternative Names}) é
bloqueado pelo \textit{Microsoft Edge} em sua configuração padrão e
gera aviso intersticial \texttt{NET::ERR\_CERT\_AUTHORITY\_INVALID} no
\textit{Google Chrome}; sua adoção pelo IRIS \textit{Mobile}
demandaria a inserção de \textit{trust anchor} dedicado ou a adoção
de \textit{certificate pinning}, cujo custo de implementação foi
considerado desproporcional ao escopo da validação.

Na mesma chave, as \textbf{credenciais e os materiais criptográficos}
hoje em uso seguem perfil declaradamente de desenvolvimento: os
arquivos \texttt{appsettings.NodeA.json} e
\texttt{appsettings.NodeB.json} embarcam diretamente os pares de
chaves RSA dos nós, os blocos PFX em base64 e suas senhas; as senhas
do PostgreSQL e do Redis seguem o padrão sugerido pelo \textit{template}
de \textit{compose}; e os arquivos \texttt{.env.production*} do IRIS
replicam os mesmos PFXs no \textit{bundle} servido pelo navegador.

A imagem do NPI é construída com
\texttt{BUILD\_CONFIGURATION=Debug}, mantida assim por simplicidade
durante o experimento. De modo análogo, o \textit{controller}
\texttt{TestingController}, que expõe rotinas auxiliares de geração de
certificados e de inspeção de canais, responde anonimamente em todos
os ambientes --- situação aceitável durante o experimento, mas
inadequada para uma operação produtiva.

Por fim, a \textbf{política de \textit{deallocate}} adotada para
conter custos da máquina virtual durante o ciclo acadêmico tem como
efeito colateral a possibilidade de troca do IP público da máquina
quando ela é reativada, o que implica reconstrução dos \textit{bundles}
do IRIS \textit{web} e do IRIS \textit{Mobile} e reemissão do
certificado autoassinado utilizado pelo NGINX.

% ----------------------------------------------------------------------
\subsection{Limitações do protocolo de sincronização}\label{subsec:limitacoesProtocolo}
% ----------------------------------------------------------------------

A análise honesta do protocolo descrito na
Subseção~\ref{subsec:NPITrocaDados} revela quatro limitações
reconhecidas que delimitam o escopo da validação experimental e
apontam para evoluções naturais do padrão.

A primeira é a \textbf{semântica de \textit{last-write-wins} sobre o
atributo \texttt{UpdatedAt}}, adotada como mecanismo de resolução de
conflitos tanto para o vocabulário SNOMED quanto para as entidades
transacionais. Essa escolha é vulnerável ao desvio de relógio entre
nós e não preserva a causalidade das modificações --- um problema
clássico nos sistemas eventualmente consistentes.

A segunda é a \textbf{ausência de \textit{checksum} explícito por
arquivo binário}. A integridade fim a fim é hoje garantida em duas
camadas: a cifra autenticada (AES-256-GCM) na camada de transporte e
a transação ACID na importação do lote no NPI consumidor. Reconhece-se,
contudo, que essa estratégia não detecta corrupção silenciosa em
casos limítrofes, como falhas de \textit{hardware} no \textit{blob
storage} de origem.

A terceira é a \textbf{sobrecarga de codificação}. O \textit{payload}
binário é serializado em Base64 dentro do envelope JSON, o que
acrescenta cerca de 33\,\% ao tamanho efetivo da transferência. A
compressão de transporte (por exemplo, GZIP) não foi habilitada na
implementação atual.

A quarta é a \textbf{descoberta manual entre nós}. O operador do
NPI consumidor registra previamente, em sua tabela local, o endereço
e o certificado esperado do NPI produtor, sem mecanismos de descoberta
distribuída como DNS-SD ou registros baseados em \textit{Distributed
Ledger Technology}. Essa decisão é deliberada --- dados clínicos
exigem federação por confiança explícita ---, mas restringe a
escalabilidade do padrão a uma rede com algumas dezenas de nós.

% ======================================================================
\section{Trabalhos futuros}\label{sec:trabalhosFuturos}
% ======================================================================

Os caminhos naturais de evolução do trabalho, organizados por
componente, são apresentados nesta seção. O agrupamento mantém
correspondência direta com as limitações sintetizadas na seção
anterior, de modo que o leitor possa percorrer cada limitação até sua
proposta de evolução.

% ----------------------------------------------------------------------
\subsection{Evoluções do PRISM e do NPI}\label{subsec:trabalhosFuturosPrism}
% ----------------------------------------------------------------------

A evolução natural da resolução de conflitos é a adoção de
\textit{vector clocks} acompanhados de uma regra determinística de
desempate por identificador de nó, em substituição ao
\textit{last-write-wins} baseado em \textit{timestamp}. Essa
mudança preserva a causalidade das modificações concorrentes e é
particularmente relevante para o domínio dos dados clínicos do
voluntário, em que edições simultâneas e independentes podem refletir
intenções clínicas distintas.

A integridade fim a fim pode ser endurecida pela inclusão de um
\textit{hash} SHA-256 por arquivo binário, calculado no NPI de origem
e revalidado no de destino. Em paralelo, a habilitação de compressão
GZIP no nível de transporte, ou a substituição da serialização em
Base64 por um canal binário paralelo, reduziria a sobrecarga de
codificação atualmente assumida pelo protocolo.

Do ponto de vista operacional, a publicação dos serviços sob um nome
de domínio mediada por DNS abriria caminho para a substituição do
certificado autoassinado por um emitido por autoridade confiável (por
exemplo, via \textit{Let's Encrypt}), eliminando o aviso intersticial
nos navegadores e dispensando o \textit{trust anchor} dedicado no
IRIS \textit{Mobile}. Os segredos hoje embarcados em arquivos de
configuração migrariam para um cofre dedicado --- como o \textit{Azure
Key Vault} --- com rotação periódica, e os contêineres receberiam
apenas referências aos seus identificadores em tempo de partida. A
imagem do NPI seria construída em \texttt{Release}, e o
\texttt{TestingController} seria bloqueado por ambiente em
implantações produtivas.

Já a abrangência da validação seria ampliada por uma topologia
multi-hospedeira na qual cada NPI residiria em sua própria máquina
virtual ou \textit{cluster}, possivelmente em regiões geograficamente
distintas, permitindo caracterizar a latência interinstitucional, as
partições de rede e a divergência de relógios em condições realistas.

Por fim, a federação seria estendida para mais de dois nós e para
modalidades de biossinal além do sEMG (eletroencefalografia,
eletrocardiografia, oximetria, acelerometria), exercendo de fato a
agnosticidade do contrato exposto pelo NPI.

% ----------------------------------------------------------------------
\subsection{Evoluções do IRIS Mobile}\label{subsec:trabalhosFuturosIrisMobile}
% ----------------------------------------------------------------------

%% NOTA: as evoluções listadas abaixo correspondem a soluções para as
%% cinco limitações descritas na Subseção \ref{subsec:limitacoesIRISMobile}.

A evolução natural do \texttt{SyncService} é uma fila de
sincronização persistente, materializada por exemplo como uma tabela
auxiliar no SQLite que registre histórico de tentativas, último erro
observado e próxima janela de envio. Essa mudança permite ao
aplicativo sobreviver a reinícios prolongados sem perder o estado da
fila e abre caminho para políticas mais sofisticadas de
\textit{backoff} adaptativo.

O suporte a múltiplos dispositivos Bluetooth simultâneos exige uma
reformulação substancial do \texttt{BluetoothContext}, com a
introdução de mecanismos de \textit{ownership} e de roteamento entre
dispositivos distintos. Cenários multicanal --- por exemplo, sEMG e
eletrocardiografia em paralelo --- emergiriam naturalmente dessa
reformulação.

A análise visual mais sofisticada do sinal demandará tratar a
\textit{performance} de re-renderização como questão de primeira
ordem, possivelmente migrando o desenho para uma camada nativa via
\textit{Skia} ou \textit{OpenGL}. Sobre essa base, é possível
introduzir filtragens em tempo real (passa-faixa, retificação,
envoltória), marcadores temporais sobrepostos ao traço e modos
auxiliares de visualização.

A granularidade das anotações pode ser expandida para um modelo com
intervalos temporais (\texttt{startedAt} e \texttt{endedAt} relativos
ao início da gravação), com possível interface de marcação direta
sobre o gráfico durante a coleta. Esse modelo viabiliza usos clínicos
que demandam precisão sub-segundo e rótulos temporais fechados.

A migração do \textit{firmware} do NeuroEstimulator para um perfil
BLE compatível com a especificação de \textit{streaming}, ou a
introdução de um \textit{gateway} intermediário que traduza SPP em
BLE/Wi-Fi, abriria caminho para uma cobertura efetivamente
multiplataforma sem reescrita das demais camadas do IRIS \textit{Mobile}.

% ----------------------------------------------------------------------
\subsection{Evoluções do NeuroEstimulator}\label{subsec:trabalhosFuturosNeuroestimulator}
% ----------------------------------------------------------------------

Caso o requisito de não modificação do \textit{hardware} seja
relaxado, a evolução natural do dispositivo é a substituição do
ADS1115 por um conversor de taxa mais alta, capaz de cobrir a banda
fisiológica completa do sEMG (até 500\,Hz) sem necessidade de
decimação agressiva. Em paralelo, a migração do módulo HC-06 para um
rádio BLE 5 ou Wi-Fi de maior largura de banda eliminaria o gargalo
de comunicação atualmente imposto pelo Bluetooth Classic e abriria a
porta para a portabilidade do aplicativo a outras plataformas.

% ----------------------------------------------------------------------
\subsection{Validação clínica e operacional ampliada}\label{subsec:trabalhosFuturosValidacao}
% ----------------------------------------------------------------------

A presente investigação valida o contrato de interoperabilidade do
PRISM e a mecânica de federação entre nós; uma etapa subsequente
natural é a validação clínica em um cenário com múltiplos voluntários,
protocolos de pesquisa diversos e profissionais de saúde como
operadores. A coleta de feedback estruturado de pesquisadores e
clínicos sobre o IRIS e o IRIS \textit{Mobile} permitiria ajustar as
decisões de \textit{User Experience} apresentadas no
Capítulo~\ref{cap:desenvolvimento}.

Em uma frente complementar, a aplicação do PRISM a outras modalidades
de biossinal --- eletroencefalografia, eletrocardiografia, oximetria
--- e a condução de um estudo multicêntrico real, com nós em
instituições distintas e enlaces geograficamente distribuídos, faria
emergir aspectos não capturados pelo recorte aqui adotado e
caracterizaria efetivamente a contribuição do padrão a um ambiente de
operação produtiva.

Por fim, um exercício de avaliação cruzada com as iniciativas
governamentais examinadas na Seção~\ref{sec:estadoDaArte} (RNDS,
USCDI, EHDS, NHS, \textit{Shared Pan-Canadian Interoperability
Roadmap} e \textit{My Health Record}) permitiria identificar pontos
de adesão semântica e divergência operacional, orientando a evolução
do PRISM no sentido da convergência com padrões consolidados.

% ======================================================================
\section{Considerações finais}\label{sec:fechamento}
% ======================================================================

O posicionamento do PRISM em relação às iniciativas examinadas no
estado da arte é deliberadamente complementar, e não competitivo. As
plataformas governamentais analisadas (RNDS, USCDI, EHDS, NHS,
\textit{Shared Pan-Canadian Interoperability Roadmap} e
\textit{My Health Record}) atuam em escala nacional e sob marcos
regulatórios robustos, com forte ênfase em prontuários eletrônicos e
em fluxos clínico-assistenciais. O PRISM, por sua vez, dirige-se ao
cenário de pesquisa clínico-hospitalar, com foco em projetos de menor
escala que necessitam compartilhar dados estruturados e biossinais
entre instituições sem incorrer no custo de adoção integral de uma
plataforma nacional.

Reconhece-se, com a honestidade que o trabalho acadêmico exige, que o
PRISM, em sua forma atual, valida o \textit{contrato} de
interoperabilidade e a \textit{mecânica} da federação entre nós, mas
não a sua \textit{eficácia clínica em produção} --- objetivo que
requer estudos multicêntricos com múltiplos voluntários e
profissionais de saúde, conforme contemplado na
Subseção~\ref{subsec:trabalhosFuturosValidacao}.

A despeito dessas limitações, o conjunto de artefatos produzido --- o
NPI, o IRIS, o IRIS \textit{Mobile} e o \textit{firmware} adaptado do
NeuroEstimulator --- é diretamente reutilizável por outros
pesquisadores que desejem adotar o PRISM em seus próprios projetos,
e os pontos de evolução identificados ao longo deste capítulo
configuram um plano de trabalho claro e documentado para a
continuidade do estudo. Espera-se, com isso, contribuir para a
redução do esforço de normalização de dados em pesquisas com
biossinais e, em consequência, para o avanço de soluções baseadas em
aprendizado de máquina que dependam de bases de dados clínicas de
qualidade.
